\SectionStart
\section{模型评价与推广}
\subsection{模型优点}
\WritingContract
  {从整篇证据链中概括本文模型为何可靠或实用。}
  {围绕具体约束表达、跨问接口、基线比较、检验覆盖和可复现性归纳优点。}
  {每项优点均能回指正文结果或检验，不引入新结果。}
  {不得评价算法的教科书优点，不得使用“先进、高效、准确”等无证据形容词。}
\subsection{模型局限}
\WritingContract
  {诚实界定结论在什么条件下可能失效。}
  {说明关键假设、数据范围、参数口径、求解近似和未覆盖场景造成的具体限制。}
  {局限与各问边界主张、风险探针或敏感性结果一致。}
  {不得写“数据有限”“时间仓促”等通用套话，不得把尚未验证的改进说成既有能力。}
\subsection{可推广方向}
\WritingContract
  {说明模型迁移到新数据、新区域或新约束时需要哪些具体改动。}
  {给出可执行的参数重估、约束扩展、数据补充或算法替换，并说明触发条件和预期解决的局限。}
  {推广方向由已识别的局限和接口设计推出，不增加未经实验验证的性能承诺。}
  {不得泛称“进一步优化模型”，不得重复各问结论，不得宣称超出证据范围的普适性。}
